\section*{작성·도해·검산과 판본 안내}
\addcontentsline{toc}{section}{작성·도해·검산과 판본 안내}

책임 편집자는 김현중이다. 이 v3는 사용자가 지정한 학습 주제와 개정 방향에 따라
편집부(AEDAUT)가 OpenAI Codex의 집필·감수 보조로 전체 설명을 새로 정리한 한국어 학습 노트다.
기존 v2의 네 칸 위상 오라클, Walsh 무늬, 기저의 구별과 P·CX 회로를 유지하면서,
관측량의 뜻, 파울리 계수의 유도, 측정 회전과 패리티, 기대값 계산을 원래의 작은 예제에 연결했다.
IBM 문서의 번역본이 아니며 참고문헌과 독립 계산을 바탕으로 구성했다.
여섯 설명도는 본문의 행렬·규칙으로 새로 그린 벡터 도해다. 외부 그림을 복제하지 않았다.

기존 v2 공개판과 그 판본의 필기 연결은 보존한다. 같은 렉처 주소에서 판본에 맞는 공책을 사용할 수 있다.
수신 원판의 집필 모델과 원저자의 별도 검산 기록은 확인되지 않았다.
이전 공개 편집의 기록과 이번 v3의 집필·검토 기록은 구분한다.
사람 독자의 이해도 시험이나 사람의 문장별 감수를 수행한 것으로 표시하지 않는다.
이번 작업에서 확인한 에이전트 식별명과 역할은 \texttt{root}(집필·편집 보조),
\texttt{ftqc\_science}(원전 대조·과학 감수), \texttt{ftqc\_pdf}(도해·PDF 제작),
\texttt{pqc\_archive\_structure}(독자 흐름·판본 연속성·발행 검수)다.

이번 v3에서는 별도 작성한 Python 표준 라이브러리 계산으로 이름을 붙인 27개 수학 항목을 검산했다.
세 상태의 기대값, 파울리 계수와 재구성, 측정 회전, 비트 순서와 패리티,
벨 상태의 상관관계, Walsh 계수, 공통 위상을 포함한 P·CX 회로,
비가환 지수 분해의 반례, 다섯 문제의 풀이와 세 큐비트 모형을 확인했다.
최대 행렬 크기는 \(8\times8\), 항등식 검산의 최대 잔차는 약 \(8.9\times10^{-16}\),
허용 기준은 \(10^{-12}\)였다. 비가환 반례에서는 두 식의 차이가 0이 아님을 따로 확인했다.
27은 검산 묶음의 수이며 실험 횟수가 아니다.
이전 판본의 검산을 재사용한 것으로 세지 않고 이번에 직접 다시 계산했다.
SDK·QPU·실물 실험과 원저자의 별도 검산 프로그램은 실행하지 않았다.
가상 측정 기록은 본문에 명시한 교육용 예다.


\AIModelDisclosure
